\chapter{LANDASAN TEORI}


Tuliskan di sini \textbf{Pedoman Penulisan Landasan Teori:} \\
Uraikan teori-teori dasar yang mendasari analisis infrastruktur jaringan komputer Anda. JANGAN mencantumkan teori di luar bidang jaringan (seperti teori aplikasi mobile computing). Teori yang relevan meliputi:
\begin{itemize}
    \item \textbf{Konsep Dasar Jaringan}: Definisi jaringan komputer, model OSI layer, protokol TCP/IP, topologi fisik/logis, dan subnetting IP.
    \item \textbf{Perangkat Keras Jaringan}: Teori mengenai Switch, Router, Access Point, kabel UTP/FO, dan NIC.
    \item \textbf{Perangkat Lunak Jaringan}: Teori mengenai sistem operasi jaringan, tools pemantauan (\textit{Ping, Traceroute, Wireshark, MRTG}), dan protokol keamanan.
\end{itemize}

